
\color{black}
\pagecolor{white}
\begin{thebibliography}{99}

\bibitem{iwaniec}
H. Iwaniec.
\textit{Topics in Classical Automorphic Forms}.
北京：高等教育出版社，2018.

\bibitem{saito-nt1}
加藤和也、黒川信重、斎藤毅.
《数论I——Fermat的梦想和类域论》.
胥鸣伟、印林生 译.
北京：高等教育出版社，2009.

\bibitem{saito-nt2}
黒川信重、栗原将人、斎藤毅.
《数论II——岩泽理论和自守形式》.
印林生、胥鸣伟 译.
北京：高等教育出版社，2009. 

\bibitem{childress}
N. Childress.
\textit{Class Field Theory}.
北京：世界图书出版公司，2014.

\bibitem{hofstadter}
D. R. Hofstadter.
\textit{Gödel, Escher, Bach: an Eternal Golden Braid}.
北京：商务印书馆，2021.

\bibitem{yessenov}
K. Yessenov.
\textit{Dirichlet's Theorem on Primes in Arithmetic Progressions}.
2006.

\bibitem{anthony}
A. Várilly.
\textit{Dirichlet's Theorem on Arithmetic Progressions}.

\bibitem{mozart.i.}
W. A. Mozart.
\textit{Clarinet Concerto in A major, K. 622, I. Allegro}.
1791.

\bibitem{mozart.ii.}
W. A. Mozart.
\textit{Clarinet Concerto in A major, K. 622, II. Adagio}.
1791.

\bibitem{mozart.iii.}
W. A. Mozart.
\textit{Clarinet Concerto in A major, K. 622, III. Rondo: Allegro}.
1791.

\bibitem{weber.i.}
C. M. von Weber.
\textit{Clarinet Concerto No. 1 in F minor, Op. 73 (J. 114), I. Allegro}.
1811.

\bibitem{weber.ii.}
C. M. von Weber.
\textit{Clarinet Concerto No. 1 in F minor, Op. 73 (J. 114), II. Adagio ma non troppo}.
1811.

\bibitem{weber.iii.}
C. M. von Weber.
\textit{Clarinet Concerto No. 1 in F minor, Op. 73 (J. 114), III. Rondo}.
1811.

\bibitem{Mendelssohn}
F. Mendelssohn.
\textit{Capriccio in e-Moll, Op. 81, Nr. 3}.
1843.

\bibitem{amy}
A. L. Schneider.
\textit{Approaches to Form in First Movements of
Clarinet Concertos from Mozart to the Twentieth
Century}.
2016.

\end{thebibliography}

{\renewcommand{\bibname}{延伸阅读}
\begin{thebibliography}{99}

\bibitem{iwaniec-kowalski}
H. Iwaniec, E. Kowalski.
\textit{Analytic Number Theory}.
北京：高等教育出版社，2019.

\bibitem{algebraic-nt}
黎景辉.
《代数数论》.
北京：高等教育出版社，2016.

\bibitem{alan}
A. Baker.
\textit{Transcendental Number Theory}.
Cambridge University Press, 1975.

\bibitem{nishioka}
K. Nishioka.
\textit{Mahler Functions and Transcendence}.
Springer, 1996.

\bibitem{nesterenko}
Yu. V. Nesterenko.
\textit{Modular Functions and Transcendence Questions}.
1996.

\bibitem{2011.14401}
T. J. Fonseca.
\textit{A Geometric Introduction to Transcendence Questions on Values of Modular Forms}.
2020.

\end{thebibliography}}
\printindex
